\chapter[An Algebraic Charge-Lattice Duality]{A Proposed Algebraic Charge-Lattice Duality}
\label{app:duality_spin_umklapp}

The main text develops the parity-constrained charge--spin labels in
Section~\ref{sec:ch21-spin-charge} and the finite CAR charge shifts and
leakage witnesses in Section~\ref{sec:ch21-umklapp}. This appendix is
reserved for a separate, more abstract question: how a proposed involution
can exchange charge-density-wave and superconducting vertex expressions in
a formal star-algebra model. The duality is supplementary to the finite CAR
argument and does not assert a finite-Fock-space equivalence.

\section{The proposed involution and its scope}

Consider the involution on charge labels
\begin{equation}
 (N_R,N_L)\longmapsto(N_R,-N_L).
 \label{eq:a09-charge-duality}
\end{equation}
In the specified formal model it is implemented on finite-support basis
vectors by a permutation. Conjugation sends $N_L$ to $-N_L$ and $Z_L(x)$ to
$Z_L(x)^*$. With a fixed ordering $R<L$ and the associated signs in the
Klein shifts, the proposed generator identities are $DF_R=F_R$ and
$DF_L=-F_L^*$. The minus sign is required for compatibility with the
adjoint relations.

Suppose also that the oscillator map negates the right generators and fixes
the left generators, and that the parameter involution is
$P^*=(u,c,-s,g^{-1},g)$. Extending these assignments to a star-algebra map
requires checking every defining relation and the star operation. On the
specified Klein products and dressed fields, the intended identities are
\begin{align}
 DK_{\mathrm{CDW}}(x)&=\zeta^xK_{\mathrm{SC}}(x)^*,&
 DK_{\mathrm{SC}}(x)&=\zeta^xK_{\mathrm{CDW}}(x)^*,\\
 DX_{\mathrm{CDW}}(P,x)&=-X_{\mathrm{SC}}(P^*,x),&
 DX_{\mathrm{SC}}(P,x)&=-X_{\mathrm{CDW}}(P^*,x).
 \label{eq:a09-duality-generators}
\end{align}
Consequently, the formal vertices are expected to satisfy
\begin{align}
 DW_{\mathrm{CDW}}(P,t,x)
   &=\zeta^xW_{\mathrm{SC}}(P^*,t,x)^*,\\
 DW_{\mathrm{SC}}(P,t,x)
   &=\zeta^xW_{\mathrm{CDW}}(P^*,t,x)^*.
 \label{eq:a09-duality-vertices}
\end{align}
The phase $\zeta^x$ comes from the source character and the shifted Klein
ordering; it cannot be discarded as an unspecified proportionality
constant.

To establish these identities, one must first verify the involution on
generators and charge basis vectors, extend it to products and
formal-series coefficients, and then apply it to the ordered
Klein--vertex products. This is a proposed algebraic construction, not a
completed proof in the current Lean Core. In particular, it does not yet
give a finite CAR map, an intertwining of finite energy budgets, or spectral
equivalence between the two physical systems. Those statements require
additional definitions and separate proofs.
